\section{System Architecture}
\label{sec:architecture}

This section describes the architecture of the MetaFlow multi-agent learning system, including the specialized client models, coordination strategies, and knowledge distillation framework.

\subsection{Overview}
\label{subsec:overview}

The MetaFlow system consists of three primary components: (1) specialized client models trained on heterogeneous data distributions, (2) coordinator modules that aggregate predictions from multiple clients, and (3) a knowledge distillation pipeline that compresses the multi-agent system into a single student model. Figure~\ref{fig:system_overview} illustrates the overall architecture.

The system addresses the fundamental challenge of collaborative inference in federated settings where individual models possess complementary expertise but no single model achieves optimal performance across the entire input distribution.

\subsection{Client Models}
\label{subsec:client_models}

We deploy two specialized client models, denoted as Client~A and Client~B, each trained on disjoint subsets of the MNIST dataset to simulate heterogeneous data distributions commonly encountered in federated learning scenarios.

\subsubsection{Client A: Low-Digit Specialist}

Client~A is trained exclusively on digits $\{0, 1, 2\}$ from the MNIST training set. This restricted domain allows the model to develop specialized representations for low-value digits. The training objective is standard cross-entropy classification:

\begin{equation}
\mathcal{L}_A = -\sum_{i=1}^{N_A} \sum_{c \in \{0,1,2\}} y_i^c \log p_A(c | x_i)
\end{equation}

where $N_A$ is the number of training samples for Client~A, $y_i^c$ is the one-hot encoded label, and $p_A(c | x_i)$ represents the predicted probability distribution from Client~A.

\subsubsection{Client B: High-Digit Specialist with Label Noise}

Client~B is trained on digits $\{3, 4, 5, 6, 7, 8, 9\}$ with an additional challenge: 20\% symmetric label noise is injected into the training data to simulate realistic data quality issues. This noise is applied by randomly reassigning labels within the valid digit range with probability $p_{\text{noise}} = 0.2$. The training objective remains:

\begin{equation}
\mathcal{L}_B = -\sum_{i=1}^{N_B} \sum_{c \in \{3,\ldots,9\}} \tilde{y}_i^c \log p_B(c | x_i)
\end{equation}

where $\tilde{y}_i^c$ denotes the potentially noisy labels. The label noise introduces a robustness challenge and simulates the imperfect data quality often encountered in real-world federated settings.

\subsubsection{Model Architecture}

Both client models share the same convolutional neural network architecture:

\begin{itemize}
    \item \textbf{Layer 1:} Conv2D(1 $\rightarrow$ 32 channels, kernel=3$\times$3) + ReLU + MaxPool(2$\times$2)
    \item \textbf{Layer 2:} Conv2D(32 $\rightarrow$ 64 channels, kernel=3$\times$3) + ReLU + MaxPool(2$\times$2)
    \item \textbf{Layer 3:} Fully connected (64$\times$5$\times$5 $\rightarrow$ 128) + ReLU + Dropout(0.5)
    \item \textbf{Layer 4:} Fully connected (128 $\rightarrow$ 10) [logits output]
\end{itemize}

The output layer produces a 10-dimensional logit vector $\mathbf{z} \in \mathbb{R}^{10}$ for all digit classes, even though each model is trained on only a subset of digits.

\subsection{Coordination Strategies}
\label{subsec:coordinators}

Given predictions from multiple client models, the coordinator must produce a unified prediction. Let $\mathbf{z}_A, \mathbf{z}_B \in \mathbb{R}^{10}$ denote the logit outputs from Client~A and Client~B, respectively. We evaluate five coordination strategies:

\subsubsection{Confidence-Based Selection}

The confidence coordinator selects the prediction from the model exhibiting the highest maximum softmax probability:

\begin{equation}
\mathbf{z}_{\text{out}} = \begin{cases}
\mathbf{z}_A & \text{if } \max(\text{softmax}(\mathbf{z}_A)) > \max(\text{softmax}(\mathbf{z}_B)) \\
\mathbf{z}_B & \text{otherwise}
\end{cases}
\end{equation}

This approach assumes that higher confidence correlates with correctness, which may not hold when models are miscalibrated.

\subsubsection{Margin-Based Selection}

The margin coordinator computes the gap between the top two predictions for each model and selects the model with the larger margin:

\begin{equation}
m_i = p_i^{(1)} - p_i^{(2)}, \quad i \in \{A, B\}
\end{equation}

where $p_i^{(1)}$ and $p_i^{(2)}$ are the highest and second-highest probabilities from model $i$. The model with larger margin $m_i$ is selected.

\subsubsection{Ensemble Averaging}

The ensemble coordinator averages the logits from both models:

\begin{equation}
\mathbf{z}_{\text{out}} = \frac{1}{2}(\mathbf{z}_A + \mathbf{z}_B)
\end{equation}

This approach leverages the wisdom of crowds but may dilute strong predictions with weak ones.

\subsubsection{Average Logits with Softmax}

Similar to ensemble averaging, but applies temperature-scaled softmax before averaging probabilities:

\begin{equation}
p_{\text{out}}(c) = \frac{1}{2}\left(\text{softmax}(\mathbf{z}_A / \tau)(c) + \text{softmax}(\mathbf{z}_B / \tau)(c)\right)
\end{equation}

where $\tau$ is the temperature parameter (default $\tau=1.0$).

\subsubsection{Agree-Then-Route Strategy}

The Agree-Then-Route coordinator implements a two-stage decision process:

\begin{enumerate}
    \item \textbf{Agreement Check:} If $\arg\max(\mathbf{z}_A) = \arg\max(\mathbf{z}_B)$, return either prediction
    \item \textbf{Disagreement Routing:} If predictions differ, invoke a learned router to select the better model
\end{enumerate}

This strategy exploits the observation that when both models agree, the prediction is likely correct (approximately 99\% accuracy on agreement cases). The routing problem is thus isolated to the more challenging disagreement cases, which constitute approximately 79--80\% of test samples.

\paragraph{Router Training}

The router is implemented as a binary classifier $r: \mathbb{R}^{10} \rightarrow \{A, B\}$ that takes as input a feature vector $\mathbf{f} \in \mathbb{R}^{10}$ and outputs a model selection decision. For each model's logits $\mathbf{z}_i$, we extract five features:

\begin{align}
f_1 &= \max(\mathbf{z}_i) \quad \text{(top logit)} \\
f_2 &= p_i^{(1)} - p_i^{(2)} \quad \text{(margin)} \\
f_3 &= -\sum_{c=0}^{9} p_i(c) \log p_i(c) \quad \text{(entropy)} \\
f_4 &= \|\mathbf{z}_i\|_2 \quad \text{(logit norm)} \\
f_5 &= p_i^{(1)} \quad \text{(confidence)}
\end{align}

The combined feature vector is $\mathbf{f} = [\mathbf{f}_A; \mathbf{f}_B] \in \mathbb{R}^{10}$, where $[\cdot; \cdot]$ denotes concatenation.

The router is trained on a probe dataset containing samples where the models disagree and exactly one is correct. The training labels are binary: $y=1$ if Client~B is correct, $y=0$ if Client~A is correct. We use binary cross-entropy loss with class balancing:

\begin{equation}
\mathcal{L}_{\text{router}} = -\frac{1}{N}\sum_{i=1}^{N} w_i \left[y_i \log \sigma(r(\mathbf{f}_i)) + (1-y_i) \log(1-\sigma(r(\mathbf{f}_i)))\right]
\end{equation}

where $\sigma(\cdot)$ is the sigmoid function and $w_i$ adjusts for class imbalance.

\paragraph{Adaptive Enabling}

The router is only enabled if it achieves validation accuracy exceeding a baseline margin heuristic by at least $\delta = 0.002$:

\begin{equation}
\text{router\_enabled} = \mathbb{I}[\text{acc}_{\text{router}}^{\text{val}} \geq \text{acc}_{\text{margin}}^{\text{val}} + \delta]
\end{equation}

This adaptive mechanism prevents overfitting and ensures the learned router provides genuine improvement over simple heuristics.

\subsection{Knowledge Distillation}
\label{subsec:distillation}

To deploy the multi-agent system efficiently, we compress the MetaFlow ensemble (Client~A + Client~B + Coordinator) into a single student model via knowledge distillation~\cite{hinton2015distilling}.

\subsubsection{Teacher Model}

The teacher is the full MetaFlow system producing soft predictions:

\begin{equation}
\mathbf{p}_{\text{teacher}} = \text{softmax}(\mathbf{z}_{\text{coordinator}} / \tau)
\end{equation}

where $\mathbf{z}_{\text{coordinator}}$ is the output from the coordinator (e.g., Agree-Then-Route) and $\tau=2.0$ is the temperature parameter.

\subsubsection{Student Model}

The student model shares the same CNN architecture as the client models but is trained from scratch using the teacher's soft predictions. The distillation loss combines Kullback-Leibler (KL) divergence and hard label cross-entropy:

\begin{equation}
\mathcal{L}_{\text{distill}} = \alpha \cdot \text{KL}(\mathbf{p}_{\text{teacher}} \| \mathbf{p}_{\text{student}}) + (1-\alpha) \cdot \mathcal{L}_{\text{CE}}(y, \mathbf{p}_{\text{student}})
\end{equation}

where $\alpha$ controls the balance between soft and hard targets. In our implementation, we use pure soft target learning ($\alpha=1.0$):

\begin{equation}
\mathcal{L}_{\text{distill}} = \tau^2 \cdot \text{KL}\left(\text{softmax}(\mathbf{z}_{\text{teacher}}/\tau) \,\|\, \text{softmax}(\mathbf{z}_{\text{student}}/\tau)\right)
\end{equation}

The $\tau^2$ scaling factor compensates for the gradient magnitude reduction caused by temperature scaling.

\subsubsection{Training Protocol}

The student is trained on a separate probe dataset of 5,000 samples using the Adam optimizer with learning rate $10^{-3}$ for 20 epochs. This probe set is disjoint from both the client training data and the router training data to prevent information leakage.

\subsection{Data Partitioning}
\label{subsec:data_partitioning}

The MNIST dataset (70,000 total samples: 60,000 training + 10,000 test) is partitioned as follows:

\begin{itemize}
    \item \textbf{Client A Training:} Digits 0--2 from the 60K training set ($\sim$18,000 samples)
    \item \textbf{Client B Training:} Digits 3--9 from the 60K training set with 20\% label noise ($\sim$42,000 samples)
    \item \textbf{Router Probe:} 5,000 samples from remaining training data (for router training and validation)
    \item \textbf{Distillation Probe:} 5,000 samples from remaining training data, disjoint from router probe
    \item \textbf{Test Set:} Full 10,000 MNIST test samples (all digits, no noise)
\end{itemize}

This partitioning ensures that coordinator and distillation components are evaluated on the full digit distribution, testing their ability to generalize beyond the specialized training domains of individual clients.

\subsection{Training Pipeline}
\label{subsec:training_pipeline}

The complete training pipeline proceeds in four stages:

\begin{enumerate}
    \item \textbf{Client Training:} Train Client~A and Client~B independently on their respective data partitions until convergence (typically 10 epochs each)
    
    \item \textbf{Router Training:} Using the trained client models, generate predictions on the router probe dataset and train the disagreement router for 25 epochs
    
    \item \textbf{Distillation:} Using the full MetaFlow system (clients + coordinator), generate soft targets on the distillation probe dataset and train the student model for 20 epochs
    
    \item \textbf{Evaluation:} Assess all coordinators and the student model on the held-out test set across metrics including overall accuracy, disagreement-only accuracy, and routing decision correctness
\end{enumerate}

All experiments use deterministic seeding (Python random seed, PyTorch CPU seed, PyTorch CUDA seed) to ensure reproducibility.
